\subsection{Dimension of the image of a morphism}
\phantomsection
\label{subsection:IV.5.4}

\begin{proposition}[5.4.1]
\phantomsection
\label{IV.5.4.1}
Let $X$ and $Y$ be locally Noetherian preschemes, and let
$f:X\to Y$ be a morphism.
\begin{enumerate}[{\rm(i)}]
\item
\phantomsection
\label{IV.5.4.1.i}
If $f$ is quasi-finite, then
\[
\dim(X)\leq\dim(\overline{f(X)})\leq\dim(Y).
\]
\item
\phantomsection
\label{IV.5.4.1.ii}
If $f$ is surjective and open (respectively, surjective and closed),
then
\[
\dim(X)\geq\dim(Y).
\]
\end{enumerate}
\end{proposition}

\begin{proof}
\begin{enumerate}[{\rm(i)}]
\item
We may replace $f$ by $f_{\mathrm{red}}$ \sref[II]{II.6.2.4}, and
therefore suppose that $X$ and $Y$ are reduced.
Let $Z$ be the reduced closed subprescheme of $Y$ whose underlying
space is $\overline{f(X)}$.
Then $f=j\circ g$, where $j:Z\to Y$ is the canonical injection and
$g:X\to Z$ is quasi-finite
\sref[I]{I.5.2.2}, \sref[II]{II.6.2.4}.
We may thus restrict to the case in which $f(X)$ is dense in $Y$.

For every $x\in X$, the ring $\sh O_x$ is then a quasi-finite
$\sh O_{f(x)}$-module \sref[II]{II.6.2.2}; consequently
\[
\mathfrak m_{f(x)}\sh O_x
\]
is an ideal of definition of $\sh O_x$ \sref[0]{0.7.4.4}.
If $A\to B$ is a local homomorphism of Noetherian local rings, with
maximal ideal $\mathfrak m$ of $A$, and if $\mathfrak mB$ is an ideal
of definition of $B$, then
\[
\dim(B)\leq\dim(A)
\]
by \sref[0]{0.16.3.10}.
Part~(i) now follows from \sref{IV.5.1.4}.

\item
By the definition of dimension \sref[0]{0.14.1.2}, it is enough to
prove the following.
For every sequence $(y_i)_{0\leq i\leq n}$ of distinct points of $Y$
such that
\[
y_i\in\overline{\{y_{i+1}\}}
\qquad(0\leq i\leq n-1),
\]
there is a sequence $(x_i)_{0\leq i\leq n}$ of points of $X$ such
that
\[
x_i\in\overline{\{x_{i+1}\}}
\qquad(0\leq i\leq n-1)
\]
and $f(x_i)=y_i$ for every $i$.

Suppose first that $f$ is surjective and open.
We construct the $x_i$ by induction on $i$.
Surjectivity gives an $x_0\in X$ such that $f(x_0)=y_0$.
Suppose the $x_i$ have been chosen for $i\leq m$ so that
$f(x_i)=y_i$ for $i\leq m$ and
$x_i\in\overline{\{x_{i+1}\}}$ for $i<m$.
Since $f$ is open and $y_{m+1}$ is a generization of $y_m$, there is
an $x_{m+1}\in f^{-1}(y_{m+1})$ that is a generization of $x_m$
\sref{IV.1.10.3}.
The induction therefore continues.

\oldpage[IV-2]{94}

Suppose now that $f$ is surjective and closed.
We construct the $x_i$ by descending induction on $i$.
Surjectivity gives an $x_n$ such that $f(x_n)=y_n$.
Suppose the $x_i$ have been chosen with the required properties for
$i>m$.
Because $f$ is closed,
\[
f\bigl(\overline{\{x_{m+1}\}}\bigr)
  =\overline{\{f(x_{m+1})\}}
  =\overline{\{y_{m+1}\}}
\]
(Bourbaki, \emph{Topologie g\'en\'erale}, chap.~I,
3\textsuperscript{rd} ed., \S~5, n\textsuperscript{o}~4, prop.~9).
There is therefore an
$x_m\in\overline{\{x_{m+1}\}}$ such that $f(x_m)=y_m$, and the
descending induction continues.
\end{enumerate}
\end{proof}

\begin{corollary}[5.4.2]
\phantomsection
\label{IV.5.4.2}
Let $X$ and $Y$ be locally Noetherian preschemes, and let
$f:X\to Y$ be a finite morphism, hence a closed morphism.
Then
\[
\dim(X)=\dim(\overline{f(X)}).
\]
If $f$ is moreover surjective, then $\dim(X)=\dim(Y)$.
\end{corollary}

\begin{remarks}[5.4.3]
\phantomsection
\label{IV.5.4.3}
\begin{enumerate}[{\rm(i)}]
\item
\phantomsection
\label{IV.5.4.3.i}
We saw in \sref{IV.4.1.2} that if $X$ and $Y$ are preschemes locally
of finite type over a field $k$, and $f$ is a $k$-morphism, the
inequality
\[
\dim(Y)\leq\dim(X)
\]
already holds when $f$ is quasi-compact and dominant.
By contrast, if one assumes only that $X$ and $Y$ are locally
Noetherian, it is possible to have $\dim(Y)>\dim(X)$ even when $f$ is
of finite type, bijective, and a local immersion.

For example, let $A$ be a discrete valuation ring, let $K$ be its
fraction field, and let $k$ be its residue field.
Put $Y=\Spec(A)$, and denote by $a$ the generic point of $Y$ and by
$b$ its closed point; thus $\kappa(a)=K$ and $\kappa(b)=k$.
Let
\[
X=\Spec(K)\amalg\Spec(k),
\]
and let $f:X\to Y$ agree on the two summands with the respective
canonical morphisms \sref[I]{I.2.4.5}.
The morphism $f$ is plainly a bijective local immersion, since
$\{a\}$ is open in $Y$.
It is also of finite type: if $\pi$ is a uniformizer of $A$, then
$K=A[\pi^{-1}]$, so $K$ is a finitely generated $A$-algebra.
Nevertheless,
\[
\dim(X)=0
\qquad\text{and}\qquad
\dim(Y)=1.
\]

\item
\phantomsection
\label{IV.5.4.3.ii}
Let $A\subset B$ be Noetherian rings, with $B$ a finite
$A$-algebra.
Corollary~\sref{IV.5.4.2} again gives
\[
\dim(B)=\dim(A)
\]
\sref[0]{0.16.1.5}.
Suppose in addition that $A$ is a Noetherian local ring.
Then $B$ is a Noetherian semilocal ring
(Bourbaki, \emph{Alg\`ebre commutative}, chap.~IV, \S~2,
n\textsuperscript{o}~5, cor.~3 of prop.~9).
If $\mathfrak n_i$ $(1\leq i\leq r)$ are the maximal ideals of $B$,
then
\begin{equation}
\dim(A)=\sup_i\dim(B_{\mathfrak n_i}).
\tag{5.4.3.1}\label{IV.5.4.3.1}
\end{equation}

Under these hypotheses one need not have
\[
\dim(B_{\mathfrak n_i})=\dim(A)
\]
for every $i$; replacing $B$ by the direct product of $B$ with the
residue field $k$ of $A$ gives an immediate example.
There are even examples in which $A$ and $B$ are integral domains of
dimension $2$ and some of the $B_{\mathfrak n_i}$ have dimension
strictly less than $\dim(A)$ \sref{IV.5.6.11}.
We shall see later, in \sref{IV.5.6.4} and \sref{IV.5.6.10}, that
this last phenomenon cannot occur when $A$ is a quotient of a
regular local ring.
\end{enumerate}
\end{remarks}
